\BourbakiExerciseGroup

\begin{enumerate}
\def\labelenumi{\arabic{enumi})}
\tightlist
\item
  \begin{enumerate}
  \def\labelenumii{\alph{enumii})}
  \tightlist
  \item
    A subgroup H of a topological group G operates properly on G according to the external law \((s, x) \rightarrow sxs^{-1}\) if and only if G is Hausdorff and H is compact (use Propositions 2 and 4).
  \end{enumerate}
\end{enumerate}

\begin{enumerate}
\def\labelenumi{\alph{enumi})}
\setcounter{enumi}{1}
\tightlist
\item
  Let G be a topological group and let H be a subgroup of G. Then the mapping \((x,y)\to xy\) of \(H\times G\) into G is proper if and only if H is quasi-compact (consider the inverse image of e under this mapping).
\end{enumerate}

\begin{enumerate}
\def\labelenumi{\arabic{enumi})}
\setcounter{enumi}{1}
\tightlist
\item
  Give an example of a compact group operating continuously (and hence not properly, cf.~Proposition 3) on a non-Hausdorff space (cf.~§ 2, Exercise 13).
\end{enumerate}

¶ 3) A topological group G operates properly on a homogeneous space G/K if and only if K is quasi-compact. (To show that the condition is sufficient, show that the canonical mapping of G onto G/K is proper and use Chapter I, § 10, no. 1, Proposition 5.)

\begin{enumerate}
\def\labelenumi{\arabic{enumi})}
\setcounter{enumi}{3}
\item
  Give an example of a Hausdorff topological group G operating properly on a Hausdorff topological space X, such that the mapping \((s, x) \rightarrow s.x\) is not closed and the canonical mapping \(X \rightarrow X/G\) is not closed {[}cf.~Exercise 1 b){]}.
\item
  A subgroup H of a topological group G operates properly on G by left translations if and only if H is closed in G.
\end{enumerate}

* 6) For each real number \(a \geqslant 1\) , put

\[
\begin{array}{l l} f _ {a} (t) = \left(t + \frac {a (a + 2)}{a + 1}, - \frac {a}{a + 1}\right) & \text {if} \quad t <   - \frac {a}{a + 1}, \\ f _ {a} (t) = (a, t) & \text {if} \quad - \frac {a}{a + 1} \leqslant t \leqslant \frac {a}{a + 1}, \\ f _ {a} (t) = \left(- t + \frac {a (a + 2)}{a + 1}, \frac {a}{a + 1}\right) & \text {if} \quad t > \frac {a}{a + 1}. \end{array}
\]

Let \(C_{a}\) denote the set of all points \(f_{a}(t)\) as t runs through R, and let \(X \subset R^{2}\) be the union of the sets \(C_{a}\) for \(a \geqslant 1\) and the lines \(D'\) and \(D''\) , where \(D'\) (resp. \(D''\) ) is the set of points \((t, -1)\) (resp. \((t, 1)\) ) for \(t \in R\) .

\begin{enumerate}
\def\labelenumi{\alph{enumi})}
\tightlist
\item
  The group \(\mathbf{R}\) operates on \(\mathbf{X}\) according to the law \((s, z) \to s.z\) such that:
\end{enumerate}

\begin{enumerate}
\def\labelenumi{(\roman{enumi})}
\item
  \(s.(t, -1) = (s + t, -1);\)
\item
  \(s.(t,1) = (t - s,1);\)
\item
  \(s.f_{a}(t) = f_{a}(s + t)\) for \(a\geqslant 1\)
\end{enumerate}

Show that R operates continuously on X and that the four conditions of Proposition 4 are satisfied, but that X/R is not Hausdorff.

\begin{enumerate}
\def\labelenumi{\alph{enumi})}
\setcounter{enumi}{1}
\tightlist
\item
  Let S be the equivalence relation on X whose classes consist of a single point z if \(z \notin D' \cup D''\) , and are of the form \((z, -z)\) if \(z \in D'\) or \(z \in D''\) . Let \(X'\) be the quotient space X/S. The relation S is compatible with the group R operating on X (§ 2, no. 4); show that R operates continuously on \(X'\) , that the four conditions of Proposition 4 are satisfied and that \(X'/R\) is Hausdorff, but that R does not operate properly on \(X'\) .
\end{enumerate}

\begin{enumerate}
\def\labelenumi{\arabic{enumi})}
\setcounter{enumi}{6}
\item
  Let G be a Hausdorff group operating continuously on a locally compact space X, and suppose that (i) for every compact subset K of X, the restriction to \(G \times K\) of the mapping \((s, x) \to s.x\) is a closed mapping of \(G \times K\) into X; (ii) for each \(x \in X\) , \(s \to s.x\) is a proper mapping of G into X. Show that G operates properly on X (use Proposition 12).
\item
  Let \(G_{0}\) be a non-discrete Hausdorff topological group in which every compact subset is finite (§ 2, Exercise 23). Let G be the group \(G_{0}\) with the discrete topology. Show that G operates continuously but not properly on \(G_{0}\) by left translation, and is such that P(K, L) is compact for all compact subsets K, L of \(G_{0}\) (Theorem 1).
\item
  Let G be a topological group operating continuously on two spaces X and X', and let \(f: \mathbf{X} \to \mathbf{X}'\) be a continuous mapping compatible with the identity mapping of G (§ 2, no. 4). Let \(\vec{f}: \mathbf{X}/\mathbf{G} \to \mathbf{X}'/\mathbf{G}\) be the continuous mapping induced by \(f\) on the quotient spaces.
\end{enumerate}

\begin{enumerate}
\def\labelenumi{\alph{enumi})}
\item
  Show that if \(f\) is open (resp. closed, proper, injective, surjective) then so is \(\overline{f}\) . (To prove that if \(f\) is proper then \(\overline{f}\) is proper, show that the inverse image of a point of \(\mathbf{X}' / \mathbf{G}\) under \(\overline{f}\) is quasi-compact.)
\item
  Suppose that G acts properly and freely on both X and X'. Show that if \(\overline{f}\) is proper, then so is f {[}observe that the restriction of f to an orbit G.x is a homeomorphism onto G. \(f(x)\) {]}.
\end{enumerate}

¶ 10) Let G be a topological group and let X, Y be two topological spaces on which G operates continuously. Then G operates continuously on X × Y according to the law (s, (x, y)) → (s.x, s.y). Let (X × Y)/G be denoted by \(X \times {}^{G}Y\).

\begin{enumerate}
\def\labelenumi{\alph{enumi})}
\item
  Show that \(X \times {}^{G}G\) (G operating by left translations on itself) is canonically homeomorphic to X. (Remark that if U is open in X, then the union of the orbits of the points of \(U \times \{e\}\) is open in \(X \times G\) ).
\item
  Let \(Y'\) be a third topological space on which G operates continuously, and let \(f: Y \to Y'\) be a continuous mapping compatible with the identity mapping of G. The mapping \(\mathrm{I}_{X} \times f\) induces a continuous mapping \(\bar{f}: X \times {}^{G}Y \to X \times {}^{G}Y'\) . If \(f\) is open (resp. proper, injective, surjective) then so is \(\bar{f}\) (Exercise 9). Give an example in which \(f\) is closed, \(G\) operates properly on \(X\) , but \(\bar{f}\) is not closed {[}take \(X = G\) , \(X' = P\) (a space consisting of a single point) and use Exercise 4{]}.
\item
  Show that if G operates properly on X, then G operates properly on \(X \times Y\) .
\end{enumerate}

\(d)\) Show that, if \(\mathbf{Y}\) is compact, the canonical mapping \(\mathbf{X} \times^{\mathbf{G}}\mathbf{Y} \to \mathbf{X}/\mathbf{G}\) is proper {[}use \(b)\) {]}.

\begin{enumerate}
\def\labelenumi{\alph{enumi})}
\setcounter{enumi}{4}
\tightlist
\item
  Suppose that G is locally compact, and let \(G'\) be the one-point compactification of G. If \(\omega\) is the point at infinity in \(G'\) , then G operates continuously on \(G'\) according to the external law given by s.t = st if \(t \in G\) , and \(s.\omega = \omega\) ; X is then homeomorphic to an open subspace of \(X \times {}^{G}G'\) (use b)). Deduce that if G operates properly on X and if in addition X/G is locally compact, then X is locally compact {[}use c) to show that \(X \times {}^{G}G'\) is Hausdorff{]}.
\end{enumerate}

\(^{*}f)\) Give an example of a locally compact group G operating continuously on a space X which is not locally compact, such that X/G consists of a single point (cf.~§ 2, Exercise 29). *

¶ 11) Let G be a topological group and let H, K be two closed subgroups of G.

\begin{enumerate}
\def\labelenumi{\alph{enumi})}
\item
  Show that the following three conditions are equivalent: (i) H × K operates properly on G according to the external law \(((h, k), s) \rightarrow hsk^{-1}\) ; (ii) H operates properly on the homogeneous space G/K; (iii) K operates properly on the homogeneous space G/H.
\item
  If G is locally compact, then the (equivalent) conditions of a) are also equivalent to the following: for every pair of compact subsets A, B of G, the intersection HA ∩ BK is compact. In particular these conditions are satisfied if one of the subgroups H, K is compact.
\item
  Suppose that G is locally compact, and show that the following conditions are equivalent: (i) for each \(x \in G\) the mapping \(h \to h.xK\) of H into G/K is proper; (ii) for each \(x \in G\) the mapping \(k \to k.xH\) of K into G/H is proper; (iii) for each \(x \in G\) and each compact subset A of G, \(Hx \cap AK\) is compact; (iv) for each \(x \in G\) and each compact subset A of G, \(Kx \cap AH\) is compact. If one of the subgroups H, K is normal or if the left and right uniformities on G coincide (§3, Exercise 3), show that these conditions imply those of a).
\end{enumerate}

\begin{enumerate}
\def\labelenumi{\arabic{enumi})}
\setcounter{enumi}{11}
\tightlist
\item
  \begin{enumerate}
  \def\labelenumii{\alph{enumii})}
  \tightlist
  \item
    Let X be a compact space and G a topological group operating continuously on X. Suppose that each orbit A with respect to G has a non-empty interior relative to its closure \(\overline{A}\) . Show that there is at least one compact orbit (consider a minimal element in the set of compact subsets of X which are stable under G).
  \end{enumerate}
\end{enumerate}

* b) Give an example of a locally compact group operating continuously on a compact space, such that none of the orbits is compact (§ 2, Exercise 29). *

¶ 13) Let \(G\) be a locally compact group and let \(D\) be a discrete subgroup of \(G\) such that the homogeneous space \(G / D\) is compact. Show that, for each \(d \in D\) , the set of all \(sds^{-1}\) , as \(s\) runs through \(G\) , is closed in \(G\) . (Prove first, with the help of Proposition 10, that there is a compact subset \(K\) of \(G\) such that \(G = K.D\) ; then use Corollary 1 to Proposition 1).

\begin{enumerate}
\def\labelenumi{\arabic{enumi})}
\setcounter{enumi}{13}
\tightlist
\item
  \begin{enumerate}
  \def\labelenumii{\alph{enumii})}
  \tightlist
  \item
    Let G be a topological group operating continuously on a topological space X, and let \(x_0\) be a point of X such that \(s.x_0 = x_0\) for all \(s \in G\) . For each neighbourhood V of \(x_0\) and each quasi-compact subset K of G, show that the set \(\bigcap_{s \in \mathbf{K}} s.V\) is a neighbourhood of \(x_0\) .
  \end{enumerate}
\end{enumerate}

\begin{enumerate}
\def\labelenumi{\alph{enumi})}
\setcounter{enumi}{1}
\tightlist
\item
  Deduce from a) that if G is a locally compact group, V a neighbourhood of e, K a compact subset of G, then the set \(\bigcap_{s\in K}sVs^{-1}\) is a neighbourhood of e.
\end{enumerate}

¶ 15) Let G be a locally compact group and let \(G_{0}\) be the identity component in G; suppose that the quotient group \(G/G_{0}\) is compact.

\begin{enumerate}
\def\labelenumi{\alph{enumi})}
\item
  Show that G is \(\sigma\) -compact (remark that \(G_{0}\) is \(\sigma\) -compact and use Proposition 10).
\item
  Let \((\mathbf{U}_{n})\) be a decreasing sequence of neighbourhoods of e in G. Show that there is a compact normal subgroup K of G, contained in the intersection of the \(U_{n}\) , and such that the identity element of G/K has a countable fundamental system of neighbourhoods. {[}There exists a compact symmetric neighbourhood \(V_{0}\) of e in G such that \(G = V_{0}G_{0}\) . Define recursively a sequence of compact symmetric neighbourhoods \(V_{n}\) of e such that \(V_{n}^{2} \subset V_{n-1} \cap U_{n}\) and \(xV_{n}x^{-1} \subset V_{n-1}\) for each \(x \in V_{0}\) (Exercise 14 b); show that the intersection K of the \(V_{n}\) answers the question{]}.
\item
  Let X be a locally compact space which has a countable base; suppose G operates continuously on X in such a way that there is no \(s \neq e\) in G for which s.x = x for all \(x \in X\) . Show that the identity element of G has a countable fundamental system of neighbourhoods. {[}Let \((W_n)\) be a countable base of the topology of X, consisting of relatively compact sets; for each pair of integers \((m, n)\) such that \(\overline{W_m} \subset W_n\) , let \(U_{mn}\) be the set of all \(s \in G\) such that \(s \cdot \overline{W_m} \subset W_n\) ; show that the \(U_{mn}\) are open in G and apply the result of b).{]}
\end{enumerate}

\begin{enumerate}
\def\labelenumi{\arabic{enumi})}
\setcounter{enumi}{15}
\tightlist
\item
  Let G be a locally compact connected group, let K be a compact normal subgroup of G and let N be a closed normal subgroup of K such that K/N has no arbitrarily small subgroups (§ 2, Exercise 30). Show that in this case N is a normal subgroup of G. (The hypothesis on K/N implies that there is a compact neighbourhood U of e in G such that \(xNx^{-1} \subset N\) for each \(x \in U\) ).
\end{enumerate}

¶ 17) Let G be a locally compact group with no arbitrarily small subgroups {[}§ 2, Exercise 30{]}.

\begin{enumerate}
\def\labelenumi{\alph{enumi})}
\item
  Show that every point of G has a countable fundamental system of neighbourhoods (Exercise 15 b)).
\item
  There is a compact neighbourhood V of e such that for any two points x, y of V, the relation \(x^{2}=y^{2}\) implies x=y. (We may suppose G not commutative. If the result were false, there would exist two sequences \((x_{n})\) , \((y_{n})\) of points of G, tending to e, such that \(x_{n}^{2}=y_{n}^{2}\) and \(x_{n}y_{n}^{-1}=a_{n}\neq e\) . Let U be a compact neighbourhood of e which contains no subgroup of G other than \(\{e\}\) , and let \(p_{n}\) be the smallest integer p\textgreater0 such that \(a_{n}^{p+1}\notin U\) : show (by passing to a subsequence if necessary) that we may suppose that \(a=\lim_{n\to\infty}a_{n}^{p_{n}}\) exists, is not equal to e and belongs to U. Show that \(a^{-1}=a\) and thus obtain a contradiction).
\item
  Let U be a compact symmetric neighbourhood of e which does not contain any subgroup of G other than \(\{e\}\) , and let V be a neighbourhood of e. Show that there exists a number \(c(V) > 0\) such that, for each pair of integers p \textgreater{} 0, q \textgreater{} 0 such that \(p \leqslant c(V)q\) , and each element \(x \in G\) such that \(x, x^{2}, \ldots, x^{q}\) belong to U, we have \(x^{p} \in V\) . {[}Argue by contradiction, by supposing that there exist two sequences of integers \((p_{n}), (q_{n})\) such that \(\lim_{n \to \infty} (p_{n}/q_{n}) = 0\) and for each n an element \(a_{n} \in G\) such that \(a_{n}^{h} \in U\) for \(1 \leqslant h \leqslant q_{n}\) , but \(a_{n}^{p_{n}} \notin V\) ; we may also suppose that the sequence \((a_{n}^{p_{n}})\) has a limit \(a \neq e\) lying in U. Show that we should then have \(a^{m} \in U\) for every integer m \textgreater{} 0 and hence arrive at a contradiction{]}.
\end{enumerate}

\begin{enumerate}
\def\labelenumi{\arabic{enumi})}
\setcounter{enumi}{17}
\tightlist
\item
  Let G be a locally compact totally disconnected topological group whose left and right uniformities coincide. Show that every neighbourhood of e in G contains a compact open normal subgroup of G (use Corollary 1 to Proposition 14 of no. 6, and Exercise 3 of § 3).
\end{enumerate}

¶ 19) Let G be a locally compact group, let \(G_{0}\) be the identity component of G, and let H be a closed subgroup of G. Show that if the homogeneous space G/H is locally connected, then \(G_{0}H\) is open in G. Let \(H_{0}\) be the identity component of H; then there is an open subgroup L of H such that \(L/H_{0}\) is compact (no. 6, Corollary 1 to Proposition 14); show that on the one hand \(LG_{0}\) is closed in G, by using Corollary 1 of Proposition 1, and on the other hand that \(LG_{0}\) is open in G, by considering the canonical image of \(G_{0}\) in G/L and using Corollary 3 of Proposition 14 and the fact that a quotient space of a locally connected space is locally connected {[}Chapter I, § 11, no. 6, Proposition 12){]}.

* 20) Let \(\varphi\) be the canonical homomorphism \(\mathbf{R} \to \mathbf{T}\) and let \(\theta\) be an irrational number. Define a group structure on the topological space \(\mathbf{R}^2 \times \mathbf{T}^2\) by the law

\[
\begin{array}{l} \left(x _ {1}, x _ {2}, t _ {1}, t _ {2}\right) \left(x _ {1} ^ {\prime}, x _ {2} ^ {\prime}, t _ {1} ^ {\prime}, t _ {2} ^ {\prime}\right) \\ = \left(x _ {1} + x _ {1} ^ {\prime}, x _ {2} + x _ {2} ^ {\prime}, t _ {1} + t _ {1} ^ {\prime} + \varphi \left(x _ {2} x _ {1} ^ {\prime}\right), t _ {2} + t _ {2} ^ {\prime} + \varphi \left(\theta x _ {2} x _ {1} ^ {\prime}\right)\right). \end{array}
\]

G thus becomes a locally compact topological group (and even a Lie group). Show that the commutator subgroup of G is not closed. *

\begin{enumerate}
\def\labelenumi{\arabic{enumi})}
\setcounter{enumi}{20}
\tightlist
\item
  Let M be a compact space endowed with a semigroup structure such that (i) the mapping \((x,y)\to xy\) of \(M\times M\) into M is continuous; (ii) for each \(a\in M\) the relation ax=ay implies x=y.
\end{enumerate}

\begin{enumerate}
\def\labelenumi{\alph{enumi})}
\item
  Show that if \(\mathbf{F}\) is a closed subset of \(\mathbf{M}\) and \(x \in \mathbf{M}\) is such that \(xF \subset F\) , then \(xF = F\) . {[}Let \(y\) be a cluster point of the sequence \((x^n)_{n \geqslant 1}\) in \(\mathbf{M}\) ; show that \(yF = \bigcap_{n \geqslant 1} x^n F\) , and deduce that \(yxF = yF\) {]}.
\item
  Deduce from a) that if, in addition, for each \(a \in M\) the relation xa = ya implies x = y, then M is a compact topological group. (Start by showing that there is an identity element e; to show that \(x \to x^{-1}\) is continuous, argue by reductio ad absurdum, considering an ultrafilter on M which converges to e).
\end{enumerate}

¶ 22) a) Let G be a Hausdorff topological group. Show that every stable subset S of G which is either compact, or open and relatively compact, is a subgroup of G {[}use Exercise 21 above, and Exercise 28 a) of § 2{]}. Deduce that a compact commutative group is a radical group {[}§ 2, Exercise 28 d){]}.

\begin{enumerate}
\def\labelenumi{\alph{enumi})}
\setcounter{enumi}{1}
\item
  Deduce from a) that every locally compact stable subset of a compact group K is a subgroup of K.
\item
  Let G be a Hausdorff topological group and let S be a stable closed subset of G with the property that there is a subset K of G which is precompact with respect to the right uniformity of G, and such that SK = G. Show that S is a subgroup of G. (Given \(s \in S\) , define recursively a sequence of elements \(x_{i} \in K\) and a sequence of elements \(s_{i} \in S\) such that \(s^{-1}x_{i} = s_{i+1}x_{i+1}\) ; observe that if i \textless{} j then \(s^{-1}x_{i}x^{-1}_{j} \in S\) , and that for every neighbourhood V of e in G there is a pair of elements \(x_{i}, x_{j}\) such that i \textless{} j and \(x_{i}x^{-1}_{j} \in V\) ).
\end{enumerate}

\begin{enumerate}
\def\labelenumi{\arabic{enumi})}
\setcounter{enumi}{22}
\tightlist
\item
  Let \(p\) be a prime number, let \((\mathbf{G}_n)\) be an infinite sequence of topological groups each identical with the discrete group \(\mathbf{Z}/(p^2)\) , and let \(\mathbf{H}_n\) be the subgroup \((p) / (p^2)\) of \(\mathbf{G}_n\) . Let \(\mathbf{G}\) be the local direct product of the \(\mathbf{G}_n\) relative to the \(\mathbf{H}_n\) ( \(\S 2\) , Exercise 26); then \(\mathbf{G}\) is a locally compact group but is not compact. Show that the continuous homomorphism \(u: x \to px\) of \(\mathbf{G}\) into itself is not a strict morphism of \(\mathbf{G}\) onto \(u(\mathbf{G})\) , and that \(u(\mathbf{G})\) is not closed in \(\mathbf{G}\) .
\end{enumerate}
% semantic-fragments: legacy-input-seam
\relax{}
